% GENERATED FILE. Edit canonical lessons, then run npm run generate:books.
% canonical-source-hash: fnv1a64:8f3e6682a8d8ceb1
% canonical-lessons: SA-C514-parivesayati, SA-C514-bhojayati, SA-C514-asnati, SA-C514-ucchistam, SA-C514-pidhanam

\chapter{Serving the Meal}
\label{ch:sa-serving-the-meal}
\hlchaptermodality{514}

\begin{chapteropening}
\textbf{By the end of this chapter:} \emph{I can serve a meal and clear up after it.}

Complete the last lesson of chapter 514: I can serve a meal and clear up after it.
\end{chapteropening}

\section[\texorpdfstring{\sk{परिवेषयति} (pariveṣayati)}{परिवेषयति (pariveṣayati)}]{\sk{परिवेषयति} (pariveṣayati) — serves food}

\label{lesson:SA-C514-parivesayati}



\begin{quote}
Before the new one: say the Sanskrit for strains, then the Sanskrit for oily.
\end{quote}

\subsection*{You'll want to know: \sk{परिवेषयति}}
\textbf{\sk{परिवेषयति}} — \emph{pariveṣayati} — \textquotedblleft{}serves food\textquotedblright{}.

\begin{quote}
\sk{माता} \sk{अतिथये} \sk{भोजनं} \sk{परिवेषयति।}
\end{quote}

\begin{quote}
\emph{mātā atithaye bhojanaṃ pariveṣayati.}
\end{quote}

\begin{quote}
Mother serves food to the guest.
\end{quote}

Past (she …): \sk{सा} \sk{ओदनं} \sk{परिवेषितवती}, \emph{sā odanaṃ pariveṣitavatī}.

\begin{cousinweb}
\textbf{\sk{परिवेषयति}} (\emph{pariveṣayati}) is built on \emph{pari}, \textquotedblleft{}around\textquotedblright{}, and the root √viṣ, \textquotedblleft{}to be active, to work\textquotedblright{}; MW gives \emph{pariviṣ} as to serve, to wait on, to offer or dress food (MW).
\end{cousinweb}

\begin{grammarlens}[title={What you've built}]
One more word for describing what you lived and felt.
\end{grammarlens}

\subsection*{Your turn}
\begin{itemize}
\raggedright
  \item \emph{Say it:} \emph{pariveṣayati}
  \item \emph{Say it:} \emph{pariveṣayati}, once more
  \item \emph{Say it:} \emph{mātā atithaye bhojanaṃ pariveṣayati.} — Mother serves food to the guest.
  \item \emph{Recall:} say the Sanskrit for strains, then the Sanskrit for oily, then say \emph{pariveṣayati} again
\end{itemize}

\begin{tcolorbox}[breakable,colback=teal!4,colframe=teal!35!black,title={Before you move on}]
What does \sk{परिवेषयति} mean? (\textquotedblleft{}Serves food\textquotedblright{}.) Say it once more.
\end{tcolorbox}
\section[\texorpdfstring{\sk{भोजयति} (bhojayati)}{भोजयति (bhojayati)}]{\sk{भोजयति} (bhojayati) — feeds}

\label{lesson:SA-C514-bhojayati}



\begin{quote}
Before the new one: say the Sanskrit for oily, then the Sanskrit for serves food.
\end{quote}

\subsection*{You'll want to know: \sk{भोजयति}}
\textbf{\sk{भोजयति}} — \emph{bhojayati} — \textquotedblleft{}feeds\textquotedblright{}.

\begin{quote}
\sk{सा} \sk{शिशुं} \sk{भोजयति।}
\end{quote}

\begin{quote}
\emph{sā śiśuṃ bhojayati.}
\end{quote}

\begin{quote}
She feeds the baby.
\end{quote}

Past (she …): \sk{सा} \sk{बालं} \sk{भोजितवती}, \emph{sā bālaṃ bhojitavatī}.

\begin{cousinweb}
\textbf{\sk{भोजयति}} (\emph{bhojayati}) comes from the causative of the root √bhuj, \textquotedblleft{}to enjoy, to eat\textquotedblright{}: to make someone eat, to feed (MW).
\end{cousinweb}

\begin{grammarlens}[title={What you've built}]
One more word for describing what you lived and felt.
\end{grammarlens}

\subsection*{Your turn}
\begin{itemize}
\raggedright
  \item \emph{Say it:} \emph{bhojayati}
  \item \emph{Say it:} \emph{bhojayati}, once more
  \item \emph{Say it:} \emph{sā śiśuṃ bhojayati.} — She feeds the baby.
  \item \emph{Recall:} say the Sanskrit for oily, then the Sanskrit for serves food, then say \emph{bhojayati} again
\end{itemize}

\begin{tcolorbox}[breakable,colback=teal!4,colframe=teal!35!black,title={Before you move on}]
What does \sk{भोजयति} mean? (\textquotedblleft{}Feeds\textquotedblright{}.) Say it once more.
\end{tcolorbox}
\section[\texorpdfstring{\sk{अश्नाति} (aśnāti)}{अश्नाति (aśnāti)}]{\sk{अश्नाति} (aśnāti) — dines}

\label{lesson:SA-C514-asnati}



\begin{quote}
Before the new one: say the Sanskrit for serves food, then the Sanskrit for feeds.
\end{quote}

\subsection*{You'll want to know: \sk{अश्नाति}}
\textbf{\sk{अश्नाति}} — \emph{aśnāti} — \textquotedblleft{}dines\textquotedblright{}.

\begin{quote}
\sk{पिता} \sk{सायं} \sk{भोजनम्} \sk{अश्नाति।}
\end{quote}

\begin{quote}
\emph{pitā sāyaṃ bhojanam aśnāti.}
\end{quote}

\begin{quote}
Father has his meal in the evening.
\end{quote}

Past (he …): \sk{सः} \sk{भोजनम्} \sk{अशितवान्}, \emph{saḥ bhojanam aśitavān}.

\begin{cousinweb}
\textbf{\sk{अश्नाति}} (\emph{aśnāti}) comes from the root √aś, \textquotedblleft{}to eat, to consume\textquotedblright{}, whose present is \emph{aśnāti}; its past participle is \emph{aśita} (MW).
\end{cousinweb}

\begin{grammarlens}[title={What you've built}]
One more word for describing what you lived and felt.
\end{grammarlens}

\subsection*{Your turn}
\begin{itemize}
\raggedright
  \item \emph{Say it:} \emph{aśnāti}
  \item \emph{Say it:} \emph{aśnāti}, once more
  \item \emph{Say it:} \emph{pitā sāyaṃ bhojanam aśnāti.} — Father has his meal in the evening.
  \item \emph{Recall:} say the Sanskrit for serves food, then the Sanskrit for feeds, then say \emph{aśnāti} again
\end{itemize}

\begin{tcolorbox}[breakable,colback=teal!4,colframe=teal!35!black,title={Before you move on}]
What does \sk{अश्नाति} mean? (\textquotedblleft{}Dines\textquotedblright{}.) Say it once more.
\end{tcolorbox}
\section[\texorpdfstring{\sk{उच्छिष्टम्} (ucchiṣṭam)}{उच्छिष्टम् (ucchiṣṭam)}]{\sk{उच्छिष्टम्} (ucchiṣṭam) — leftovers}

\label{lesson:SA-C514-ucchistam}



\begin{quote}
Before the new one: say the Sanskrit for feeds, then the Sanskrit for dines.
\end{quote}

\subsection*{You'll want to know: \sk{उच्छिष्टम्}}
\textbf{\sk{उच्छिष्टम्}} — \emph{ucchiṣṭam} — \textquotedblleft{}leftovers\textquotedblright{}.

\begin{quote}
\sk{उच्छिष्टं} \sk{धेनवे} \sk{ददातु।}
\end{quote}

\begin{quote}
\emph{ucchiṣṭaṃ dhenave dadātu.}
\end{quote}

\begin{quote}
Please give the leftovers to the cow.
\end{quote}

A neuter noun.

\begin{cousinweb}
\textbf{\sk{उच्छिष्टम्}} (\emph{ucchiṣṭam}) is built on \emph{ud}, \textquotedblleft{}out, over\textquotedblright{}, and the root √śiṣ, \textquotedblleft{}to leave\textquotedblright{}, as its past participle; MW gives \emph{ucchiṣṭa} as left over, and the remnants of food (MW).
\end{cousinweb}

\begin{grammarlens}[title={What you've built}]
One more word for describing what you lived and felt.
\end{grammarlens}

\subsection*{Your turn}
\begin{itemize}
\raggedright
  \item \emph{Say it:} \emph{ucchiṣṭam}
  \item \emph{Say it:} \emph{ucchiṣṭam}, once more
  \item \emph{Say it:} \emph{ucchiṣṭaṃ dhenave dadātu.} — Please give the leftovers to the cow.
  \item \emph{Recall:} say the Sanskrit for feeds, then the Sanskrit for dines, then say \emph{ucchiṣṭam} again
\end{itemize}

\begin{tcolorbox}[breakable,colback=teal!4,colframe=teal!35!black,title={Before you move on}]
What does \sk{उच्छिष्टम्} mean? (\textquotedblleft{}Leftovers\textquotedblright{}.) Say it once more.
\end{tcolorbox}
\section[\texorpdfstring{\sk{पिधानम्} (pidhānam)}{पिधानम् (pidhānam)}]{\sk{पिधानम्} (pidhānam) — a lid}

\label{lesson:SA-C514-pidhanam}



\begin{quote}
Before the new one: say the Sanskrit for dines, then the Sanskrit for leftovers.
\end{quote}

\subsection*{You'll want to know: \sk{पिधानम्}}
\textbf{\sk{पिधानम्}} — \emph{pidhānam} — \textquotedblleft{}a lid\textquotedblright{}.

\begin{quote}
\sk{पात्रस्य} \sk{पिधानं} \sk{कुत्र} \sk{अस्ति}?
\end{quote}

\begin{quote}
\emph{pātrasya pidhānaṃ kutra asti?}
\end{quote}

\begin{quote}
Where is the lid of the pot?
\end{quote}

A neuter noun.

\begin{cousinweb}
\textbf{\sk{पिधानम्}} (\emph{pidhānam}) is built on \emph{pi}, a short form of \emph{api}, \textquotedblleft{}on\textquotedblright{}, and the root √dhā, \textquotedblleft{}to put\textquotedblright{}; MW gives \emph{pidhāna} as a covering, and a cover or lid (MW).
\end{cousinweb}

\begin{grammarlens}[title={What you've built}]
One more word for describing what you lived and felt.
\end{grammarlens}

\subsection*{Your turn}
\begin{itemize}
\raggedright
  \item \emph{Say it:} \emph{pidhānam}
  \item \emph{Say it:} \emph{pidhānam}, once more
  \item \emph{Say it:} \emph{pātrasya pidhānaṃ kutra asti?} — Where is the lid of the pot?
  \item \emph{Recall:} say the Sanskrit for dines, then the Sanskrit for leftovers, then say \emph{pidhānam} again
\end{itemize}

\begin{tcolorbox}[breakable,colback=teal!4,colframe=teal!35!black,title={Before you move on}]
What does \sk{पिधानम्} mean? (\textquotedblleft{}A lid\textquotedblright{}.) Say it once more.
\end{tcolorbox}
